\chapter{基于对称密钥的实例化协议与综合性能评估}
\echapter{Instantiated Symmetric-Key Based Protocols and Comprehensive Performance Evaluation}

\section{引言}
\esection{Introduction}

在第五章中，本文构建了一个基于对称密钥技术的多方隐私集合运算（MPSO）统一框架。该框架通过首先引入规范集合谓词公式（Canonical Set Predicate Formula, CSPF）对目标集合公式进行统一表示，然后提出谓词零分享（Predicative Zero-Sharing）这一新型密码学原语及其子类 —— 成员零分享（Membership Zero-Sharing）作为核心组件，成功实现了计算任意集合公式的通用 MPSO 功能，并打破了以往研究中的“功能孤岛”，解决了现有协议功能局限和缺乏统一性的问题。

然而，在安全多方计算（MPC）领域，协议的“通用性”往往意味着在处理某些具体场景时可能存在性能冗余。例如，在上述通用 MPSO 框架的设计中，为了消除按照 CSPF 中子公式的顺序生成秘密分享的位置与重构出的元素所属局部子集和在布谷鸟哈希表中的位置的之间的关联，以防止秘密分享重构阶段造成信息泄露，必须强制要求所有参与方在重构前执行一个多方秘密分享洗牌协议。然而，在某些特殊情况下，例如 MPSI 协议中，由于交集必然是 Leader 输入集合的子集，因此重构出的秘密的位置完全由其输入集合决定，这对其而言不再构成额外的信息泄露。基于这一子集属性，可以直接在 MPSI 协议中省去多方秘密分享洗牌这一步骤，优化实例化出的 MPSI 协议的效率。

为了在实际应用中实现极致的效率，本章针对 MPSO 中的典型功能，在第五章提出的通用 MPSO 框架下实例化的各具体协议进行深度的协议级优化，通过\textbf{结合特定集合公式的结构特性，裁剪通用流程中的冗余步骤}，给出了高效的 MPSI、MPCSI 和 MPSU 协议，其中的 MPSI 协议是首个在标准半诚实模型下实现\textbf{最优渐进复杂度}（即达到与不考虑隐私的明文传输方案相同量级的复杂度 —— Leader 的计算与通信复杂度为 $O(mn)$，Client 的计算与通信复杂度为 $O(n)$）的基于对称密钥技术的构造；MPSI-card 和 MPSI-card-sum 协议是首个在标准半诚实模型下实现\textbf{最优渐进复杂度}的 MPSI-card 和 MPSI-card-sum 协议；Circuit-MPSI 协议是首个在\textbf{不诚实大多数（Dishonest Majority）}设定下安全的Circuit-MPSI 构造；MPSU 协议是标准半诚实模型下基于对称密钥的另一 MPSU 方案，相比第三章提出的 MPSU 协议，在理论上实现了\textbf{更优的渐进复杂度}。

为了验证第五章提出的通用 MPSO 框架的实际可用性，我们使用 C++ 语言系统实现了述该框架下实例化出的具体协议，并在相同的硬件平台和网络环境（涵盖局域网 LAN 和广域网 WAN）下，对所有协议，包括第三章和第四章提出的 MPSU 协议与当前标准半诚实模型下各个具体子领域最先进（State-of-the-Art, SOTA）的方案进行了统一的基准测试（Benchmark）与性能对比。我们的实验数据充分证明，该框架在保证了功能通用性的同时，其实际在线效率已达到甚至超越了现有的各 SOTA 协议。

本章后续结构如下：6.2 节和 6.3 节首先介绍框架底层组件成员零分享的两个特例——全成员零分享和全非成员零分享的具体构造；6.4 节给出针对 MPSI 及其变体（MPSI, MPSI-card, MPSI-card-sum, Circuit-MPSI）的框架实例化协议的优化版本；6.5 节探讨针对 MPSU 及其变体（MPSU, MPSU-card, Circuit-MPSU）的框架实例化协议的优化版本；6.6 节 详细分析了各类协议的复杂度并给出了各个具体子领域的相关工作的全面理论对比；6.7 节展示各类协议的具体实现细节及综合性能评估结果；最后，6.8 节对本章内容进行总结。

\section{全成员零分享及其变体} 
\esection{Full Membership Zero-Sharing and Variant}

\subsection{全成员零分享}

批量全成员零分享（Batch Full Membership Zero-Sharing）是成员零分享在特定谓词下的一个特例，即当关联的公式 $Q$ 是 $m-1$ 个集合成员谓词的合取时（即 $\bigwedge_{j \in \{1,\cdots,m\} \setminus \{\mathsf{pivot}\}} x \in X_j$）。该功能专门为第 \ref{sec:opt-mpsi} 节中优化的 MPSI 及其变体设计。理想功能 $\FuncbpMZS$ 在图 \ref{fig:func-bpmzs} 中给出，其完整协议构造在图 \ref{fig:proto-bpmzs} 中描述。

\begin{figure}[!hbtp]
\begin{framed}
\begin{minipage}[center]{\textwidth}
\begin{trivlist}
\item \textbf{参数：} $m$ 个参与方 $P_1, \cdots P_m$，其中仅枢纽方 $P_\mathsf{pivot}$ 持有 $n$ 个独立元素作为输入，其余参与方持有大小为 $n$ 的集合。批量大小 $n$。有限域 $\mathbb{F}$。 
\item \textbf{功能：} 接收来自 $P_\mathsf{pivot}$ 的输入 $\textbf{x} = (x_1, \cdots, x_n)$ 以及来自每个 $P_j$ ($j \in \{1,\cdots m\} \setminus \{\mathsf{pivot}\}$) 的输入 $\textbf{X}_j = (X_{j,1}, \cdots, X_{j,n})$。对于 $1 \le i \le m$，采样 $\textbf{s}_i = (s_{i,1}, \cdots, s_{i,n}) \gets \mathbb{F}^n$，且满足：对于 $1 \le d \le n$，如果 $\bigwedge_{j \in \{1,\cdots,m\} \setminus \{\mathsf{pivot}\}} (x_d \in X_{j,d}) = 1$，则 $\sum_{1 \le i \le m} s_{i,d} = 0$。向 $P_i$ 发送 $\textbf{s}_i$。
\end{trivlist}
\end{minipage}
\end{framed}
\caption{批量全成员零分享功能 $\FuncbpMZS$}\label{fig:func-bpmzs}
\end{figure}

\begin{figure}[!hbtp]
\begin{framed}
\begin{minipage}[center]{\textwidth}
\begin{trivlist}
\item \textbf{参数：} $m$ 个参与方 $P_1, \cdots P_m$。批量大小 $n$。有限域 $\mathbb{F}$。离线阶段生成的 $n$ 个 Beaver 三元组 $([\textbf{a}], [\textbf{b}], [\textbf{c}])$，其中 $[\textbf{a}] = ([a_1], \cdots, [a_n])$，$[\textbf{b}] = ([b_1], \cdots, [b_n])$，$[\textbf{c}] = ([c_1], \cdots, [c_n])$ 且满足 $c_i = a_i \cdot b_i$。
\item \textbf{输入：} $P_\mathsf{pivot}$ 输入向量 $\textbf{x} = (x_1, \cdots, x_n)$。对于 $j \in \{1,\cdots m\} \setminus \{\mathsf{pivot}\}$，$P_j$ 输入 $\textbf{X}_j = (X_{j,1}, \cdots, X_{j,n})$。
\item \textbf{协议流程：}
\begin{enumerate}[itemsep=2pt,topsep=0pt,parsep=0pt,leftmargin=15pt]
\item 对于第 $i$ 个实例 ($1 \le i \le n$)，$P_j$ 采样 $r_{j,i}$ 并设置键集合 $K_{j,i} = X_{j,i}$，值集合 $V_{j,i} = \{-r_{j,i}, \cdots, -r_{j,i}\}$，其中 $\lvert K_{j,i} \rvert = \lvert V_{j,i} \rvert$。
\item $P_\mathsf{pivot}$ 和 $P_j$ 调用 $\FuncbOPPRF$。其中 $P_j$ 作为发送方 $\Sd$ 输入 $(K_{j,1},\cdots,K_{j,n})$ 和 $(V_{j,1},\cdots,V_{j,n})$；$P_\mathsf{pivot}$ 作为接收方 $\Rcv$ 输入 $\textbf{x}$ 并接收输出 $\textbf{u}_{j}$。
\item $P_\mathsf{pivot}$ 设置其秘密分享份额为 $\textbf{r}_{\mathsf{pivot}} = \sum_{j\in \{1,\cdots m\} \setminus \{\mathsf{pivot}\}} \textbf{u}_{j}$。$P_j$ 设置其份额为 $\textbf{r}_{j} = (r_{j,1}, \cdots, r_{j,n})$。此时，所有参与方持有一个包含 $n$ 个秘密分享的向量 $[\textbf{r}] = ([r_1], \cdots, [r_n])$。
\item 所有参与方使用 Beaver 三元组 $([\textbf{a}], [\textbf{b}], [\textbf{c}])$ 执行 $n$ 次安全乘法计算，得到 $[\textbf{s}]$，即 $[s_i] = [r_i] \cdot [b_i]$ ($1 \le i \le n$)。
\end{enumerate}
\end{trivlist}
\end{minipage}
\end{framed}
\caption{批量全成员零分享协议 $\ProtobpMZS$}\label{fig:proto-bpmzs}
\end{figure}

\begin{trivlist}
\item \textbf{复杂度分析.} 在批量全成员零分享协议（图~\ref{fig:proto-bpmzs}）中，各阶段的开销计算如下：
\begin{itemize}
  \item $P_\mathsf{pivot}$ 与每个 $P_j$ 执行批量大小为 $n$ 的 $\FuncbOPPRF$。结合哈希分桶技术时，每组集合大小 $\lvert X_{j,i} \rvert = O(1)$。我们遵循相关文献利用向量不经意线性求值（Vector Oblivious Linear Evaluation, VOLE）和不经意键值存储（OKVS）来实例化 $\FuncbOPPRF$。利用通信均摊技术，计算 $n$ 个实例的总通信量与总元素数量 $3n$ 相当。这使得批量 OPPRF 的计算复杂度随 $n$ 线性扩展。因此，在此阶段，$P_\mathsf{pivot}$ 的计算和通信复杂度均为 $O(mn)$，每个 $P_j$ 的计算和通信复杂度均为 $O(n)$。
  \item 参与方利用 Beaver 三元组执行 $n$ 次安全乘法，并指定 $P_\mathsf{pivot}$ 作为领导者。这需要 $P_\mathsf{pivot}$ 产生 $O(mn)$ 的开销，每个 $P_j$ 产生 $O(n)$ 的开销。
\end{itemize}
综上所述，在在线阶段，$P_\mathsf{pivot}$ 的整体计算和通信复杂度为 $O(mn)$，每个 $P_j$ 的计算和通信复杂度为 $O(n)$。
\end{trivlist}

\subsection{带载荷的全成员零分享}

为了支持 MPSI-card-sum 功能，我们引入了一种称为“带载荷的批量全成员零分享（Batch Full Membership Zero-Sharing with Payloads）”的变体。在此变体中，$P_\mathsf{pivot}$ 持有元素 $x$，其他参与方持有元素集合及关联的载荷（Payload）集合。如果公式为真（即 $x$ 属于所有集合），参与方将获得 0 的秘密分享，以及 $x$ 的所有关联载荷之和的秘密分享；否则获得两个随机值的分享。

其构造思想与图 \ref{fig:proto-bpmzs} 类似，核心在于将载荷集合编码到 OPPRF 发送方的输入值集合中。具体而言，$P_j$ 采样两个份额 $r, w$，并设置值集合为一对元组 $(-r, v_i - w)$，其中 $v_i$ 是元素的关联载荷。通过 $\FuncbOPPRF$，$P_\mathsf{pivot}$ 最终可以获得 $(r', w')$。如果 $x$ 匹配，则 $(r', w') = (-r, v - w)$。该变体理想功能 $\FuncbpMZSp$ 的定义详见图 \ref{fig:func-mzsp}，协议细节见图 \ref{fig:proto-bpmzsp}。其渐进复杂度与不带载荷的版本保持一致。

\begin{figure}[!hbtp]
\begin{framed}
\begin{minipage}[center]{\textwidth}
\begin{trivlist}
\item \textbf{参数：} $m$ 个参与方 $P_1, \cdots P_m$。批量大小 $n$。有限域 $\mathbb{F}$ 及载荷域 $\mathbb{F'}$。映射函数 $\mathsf{payload}_j()$ 用于关联元素与载荷。
\item \textbf{功能：} 接收 $P_\mathsf{pivot}$ 的输入 $\textbf{x}$，以及 $P_j$ 的集合 $\textbf{X}_j$ 和载荷 $\textbf{V}_j$。对于 $1 \le i \le m$，采样 $\textbf{s}_i \gets \mathbb{F}^n, \textbf{w}_i \gets \mathbb{F'}^n$，满足：对于 $1 \le d \le n$，如果 $\bigwedge (x_d \in X_{j,d}) = 1$，则 $\sum s_{i,d} = 0$ 且 $\sum w_{i,d} = \sum v_{j,d}$（其中 $v_{j,d} = \mathsf{payload}_j(x_d)$）。向 $P_i$ 输出 $(\textbf{s}_i, \textbf{w}_i)$。
\end{trivlist}
\end{minipage}
\end{framed}
\caption{带载荷的批量全成员零分享功能 $\FuncbpMZSp$}\label{fig:func-mzsp}
\end{figure}

\begin{figure}[!hbtp]
\begin{framed}
\begin{minipage}[center]{\textwidth}
\begin{trivlist}
\item \textbf{参数：} 同上，另需 $n$ 个 Beaver 三元组 $([\textbf{a}], [\textbf{b}], [\textbf{c}])$。
\item \textbf{协议流程：}
\begin{enumerate}[itemsep=2pt,topsep=0pt,parsep=0pt,leftmargin=15pt]
\item 对于第 $i$ 个实例，$P_j$ 采样 $(r_{j,i}, w_{j,i})$。设置键集合 $K_{j,i} = X_{j,i} = (x_{j,i,1}, \dots)$，值集合 $V_{j,i}' = \{(-r_{j,i}, v_{j,i,1} - w_{j,i}), \dots\}$。
\item $P_\mathsf{pivot}$ 和 $P_j$ 调用 $\FuncbOPPRF$。$P_j$ 作为 $\Sd$ 输入键值对，$P_\mathsf{pivot}$ 作为 $\Rcv$ 接收 $(\textbf{r}_{j}',\textbf{w}_{j}')$。
\item $P_\mathsf{pivot}$ 设置份额 $\textbf{r}_{\mathsf{pivot}} = \sum \textbf{r}_{j}'$，$\textbf{w}_{\mathsf{pivot}} = \sum \textbf{w}_{j}'$。$P_j$ 设置份额 $\textbf{r}_{j}$ 和 $\textbf{w}_{j}$。
\item 所有参与方使用三元组计算 $[\textbf{s}] = [\textbf{r}] \cdot [\textbf{b}]$。
\end{enumerate}
\end{trivlist}
\end{minipage}
\end{framed}
\caption{带载荷的批量全成员零分享协议 $\ProtobpMZSp$}\label{fig:proto-bpmzsp}
\end{figure}

\section{全非成员零分享} 
\esection{Full Non-Membership Zero-Sharing}

批量全非成员零分享（Batch Full Non-Membership Zero-Sharing）是当关联公式 $Q$ 是 $m-1$ 个集合非成员谓词的合取时（即 $\bigwedge_{j \in \{1,\cdots,m\} \setminus \{\mathsf{pivot}\}} x \notin X_j$）的特例。该功能专门为第 \ref{sec:opt-mpsu} 节中的优化的 MPSU 及其变体设计。理想功能 $\FuncbpNMZS$ 在图 \ref{fig:func-bpnmzs} 中给出，协议构造见图 \ref{fig:proto-bpnmzs}。

该协议的构造核心在于将底层的秘密分享隐私成员测试（ssPMT，输出布尔秘密分享）高效地转化为算术域 $\mathbb{F}$ 上的零分享。我们利用双选择比特的 ROT 协议完成了这一转化，实现了线性开销。

\begin{figure}[!hbtp]
\begin{framed}
\begin{minipage}[center]{\textwidth}
\begin{trivlist}
\item \textbf{参数：} $m$ 个参与方 $P_1, \cdots P_m$。批量大小 $n$。有限域 $\mathbb{F}$。 
\item \textbf{功能：} 接收输入 $\textbf{x}$ 和 $\textbf{X}_j$。对于 $1 \le i \le m$，采样 $\textbf{s}_i \gets \mathbb{F}^n$，且满足：如果 $\bigwedge_{j \in \{1,\cdots,m\} \setminus \{\mathsf{pivot}\}} x_d \notin X_{j,d} = 1$，则 $\sum_{1 \le i \le m} s_{i,d} = 0$。向 $P_i$ 发送 $\textbf{s}_i$。
\end{trivlist}
\end{minipage}
\end{framed}
\caption{批量全非成员零分享功能 $\FuncbpNMZS$}\label{fig:func-bpnmzs}
\end{figure}

\begin{figure}[!hbtp]
\begin{framed}
\begin{minipage}[center]{\textwidth}
\begin{trivlist}
\item \textbf{参数：} $m$ 个参与方，批量大小 $n$。Beaver 三元组 $([\textbf{a}], [\textbf{b}], [\textbf{c}])$。
\item \textbf{协议流程：}
\begin{enumerate}[itemsep=2pt,topsep=0pt,parsep=0pt,leftmargin=15pt]
\item $P_\mathsf{pivot}$ 和 $P_j$ 调用 $\FuncbssPMT$。对于第 $i$ 个实例，$P_j$ 输入 $X_{j,i}$ 获得布尔份额 $e_{j,i}^0$，$P_\mathsf{pivot}$ 输入 $x_i$ 获得布尔份额 $e_{j,i}^1$。 
\item $P_\mathsf{pivot}$ 和 $P_j$ 调用 $n$ 个 ROT 实例。$P_j$ 作为发送方 $\Sd$ 接收随机数 $r_{j,i}^0,r_{j,i}^1$，$P_\mathsf{pivot}$ 作为接收方 $\Rcv$ 输入选择比特 $e_{j,i}^1$ 并接收 $r_{j,i}^{e_{j,i}^1}$。$P_\mathsf{pivot}$ 设置 $\textbf{r}'_{j} = (r_{j,1}^{e_{j,1}^1},\cdots,r_{j,n}^{e_{j,n}^1})$。 
\item $P_\mathsf{pivot}$ 设置份额 $\textbf{r}_{\mathsf{pivot}} = \sum_{j} \textbf{r}'_{j}$。$P_j$ 设置份额 $\textbf{r}_{j} = (-r_{j,1}^{e_{j,1}^0}, \cdots, -r_{j,n}^{e_{j,n}^0})$。
\item 所有参与方使用三元组计算 $[\textbf{s}] = [\textbf{r}] \cdot [\textbf{b}]$。
\end{enumerate}
\end{trivlist}
\end{minipage}
\end{framed}
\caption{批量全非成员零分享协议 $\ProtobpNMZS$}\label{fig:proto-bpnmzs}
\end{figure}

\begin{trivlist}
\item \textbf{复杂度分析.} 在协议中，$P_\mathsf{pivot}$ 与每个 $P_j$ 执行 $\FuncbssPMT$。基于本文第三章的结论，该操作具有线性的计算和通信复杂度。后续的 ROT 操作也可在离线阶段预处理，使得在线阶段 $P_\mathsf{pivot}$ 仅需发送掩码后的选择比特。最后的三元组乘法同样保持线性开销。因此，在线阶段 $P_\mathsf{pivot}$ 的复杂度为 $O(mn)$，每个 Client $P_j$ 为 $O(n)$。
\end{trivlist}

\section{多方隐私集合求交与交集计算} 
\esection{Multi-Party Private Set Intersection and Computation on Set Intersection}

考虑计算目标为一个可构造交集 $Y$，其规范集合谓词公式 $\psi$ 相对于 $X_1$ 是局部的。例如 $X_1 \cap X_2 \cap X_3$ 相对于 $X_1$ 的局部剩余公式为 $x \in X_2 \land x \in X_3$。在这种情况下，$Y$ 的所有元素必然属于 $P_1$（即 $Y \subseteq X_1$）。这意味着秘密分享的排列顺序完全由 $X_1$ 决定。对于 MPSI 而言，由于 $P_1$ 最终将获得结果集，这种原本会被视为泄露的“顺序信息”对 $P_1$ 而言不再构成真实的隐私泄露。

因此，在 MPSI 协议中，我们可以\textbf{省去多方秘密分享洗牌（Shuffle）步骤}。当 $P_1$（作为枢纽方 $P_{\mathsf{pivot}}$）与其他各方调用完批量全成员零分享后，各方可以直接向 $P_1$ 重构秘密。如果重构的秘密为 0，则说明 $P_1$ 对应的元素属于交集。这一优化极大地提升了在线效率。
（注意：对于 MPSI-card，为了防止 $P_1$ 知道具体哪些元素在交集内，仍然需要执行洗牌操作）。

此外，在这种针对性优化下，协议开销不再依赖于元素的比特长度 $l$，因为参与方可以提前在本地对元素进行哈希压缩。为了保证正确性，必须确保哈希过程在 $P_1$ 和其他参与方的元素之间不引入冲突，因此哈希函数的输出长度至少为 $\sigma + \log_2 (m-1) + 2 \log_2 n$。
% 例如 CSPF 表示中只包含一个子公式并且该子公式关于 Leader 的输入集合是局部的，这时目标集合划分中只包含一个局部子集并且秘密分享的顺序由 Leader 本身的输入集合决定，因此重构出秘密的位置并不会泄露给 Leader 任何额外的信息， 各方根本不需要执行多方秘密分享洗牌协议。

\subsection{多方隐私集合求交、交集求势和基于电路的MPSI}

MPSI、MPSI-card 以及 Circuit-MPSI 的完整优化协议描述在图 \ref{fig:proto-mpsi} 中给出。
值得注意的是，现有的最先进 MPSI 协议 \cite{WuYC24,ChenDGB22} 通常需要参与方之间两两在线交互（Pairwise Interaction），这导致 Client 的复杂度依赖于参与方数量或腐化阈值 $t$。而我们的协议引入了\textbf{星型通信拓扑（Star Topology）}，使得 Leader 的计算和通信复杂度为 $O(mn)$，每个 Client 仅为 $O(n)$。这达到了与不安全的朴素哈希方案相同的最优复杂度。据我们所知，这是首次在标准半诚实模型下同时实现 MPSI 和 MPSI-card 最优复杂度的构造。

\begin{figure}[!hbtp]
\begin{framed}
\begin{minipage}[center]{\textwidth}
\textbf{参数：} $m$ 个参与方。集合大小 $n$。有限域 $\mathbb{F}$。哈希参数 $B$。

\textbf{输入：} $P_i$ 输入集合 $X_i$。

\begin{enumerate}[itemsep=2pt,topsep=0pt,parsep=0pt,leftmargin=15pt]
\item \textbf{哈希分桶.} $P_1$ 使用布谷鸟哈希 $\mathcal{C}_1^b \gets \Cuckoo(X_1)$。$P_j$ ($j>1$) 使用简单哈希 $\mathcal{T}_j^b \gets \Simple(X_j)$。
\item \textbf{批量全成员零分享.} 所有参与方调用 $\FuncbpMZS$，$P_1$ 作为 $P_\mathsf{pivot}$ 输入 $\mathcal{C}_1$。$P_i$ 获得份额 $\textbf{s}_i$。
\end{enumerate}

后续操作依功能而定：

    \textbf{MPSI.} 
    \begin{enumerate}[itemsep=2pt,topsep=0pt,parsep=0pt,leftmargin=15pt] \addtocounter{enumi}{2}
        \item $P_j$ 发送 $\textbf{s}_j$ 给 $P_1$。$P_1$ 重构秘密 $\textbf{s} = \sum \textbf{s}_i$。对于每个 $b$，若 $s_b = 0$，则将 $\mathcal{C}_1^b$ 中的元素加入输出结果 $Y$ 中。
    \end{enumerate}
    \textbf{MPSI-card.} 
    \begin{enumerate}[itemsep=2pt,topsep=0pt,parsep=0pt,leftmargin=15pt] \addtocounter{enumi}{2}
        \item $P_i$ 将 $\textbf{s}_i$ 输入 $\FuncMS$ 获得打乱后的 $\textbf{s}_{i}'$。
        \item $P_j$ 发送 $\textbf{s}_j'$ 给 $P_1$。$P_1$ 统计 $\sum \textbf{s}_i'$ 中 0 的数量作为基数输出。
    \end{enumerate}
    \textbf{Circuit-MPSI.}
    \begin{enumerate}[itemsep=2pt,topsep=0pt,parsep=0pt,leftmargin=15pt] \addtocounter{enumi}{2}
        \item 参与方将分享份额 $s_{i,b}$ 输入到通用 MPC 电路中，电路在内部进行相等性判断与后续逻辑运算。
    \end{enumerate}
\end{minipage}
\end{framed}
\caption{MPSI / MPSI-card / Circuit-MPSI 优化协议}\label{fig:proto-mpsi}
\end{figure} 

\subsection{多方隐私交集求势与和} 

MPSI-card-sum 协议通过利用 6.2 节中“带载荷的批量全成员零分享”原语实现。在此过程中，各方获得指示交集的零分享以及载荷的秘密分享。随后经过两次洗牌操作（分别对零分享和载荷分享进行打乱），$P_1$ 统计零分享计算基数，并基于零分享构建的指示向量汇总载荷求和。其渐进复杂度同样保持在最优水平（即与 MPSI 一致）。该协议的完整流程在图 \ref{fig:proto-mpsi-sum} 中给出。

\begin{figure}[!hbtp]
\begin{framed}
\begin{minipage}[center]{\textwidth}
\begin{trivlist}
\item \textbf{参数：} $m$ 个参与方 $P_1, \cdots, P_m$。集合大小 $n$。元素长度 $l$。有限域 $\mathbb{F}$。布谷鸟哈希参数：哈希函数 $h_1, h_2, h_3$ 及桶数量 $B$。
\item \textbf{输入：} 每个参与方 $P_i$ 拥有输入集合 $X_i = \{x_i^1,\cdots, x_i^n\} \subseteq \{0,1\}^l$ 及载荷集合 $V_i = \{v_i^1,\cdots, v_i^n\}$，并具有从元素到关联载荷的映射函数 $\mathsf{payload}_i()$。

\begin{enumerate}[itemsep=2pt,topsep=0pt,parsep=0pt,leftmargin=15pt]
    \item \textbf{哈希分桶.} $P_1$ 执行布谷鸟哈希 $\mathcal{C}_1^1, \cdots, \mathcal{C}_1^B \gets \mathsf{Cuckoo}_{h_1,h_2,h_3}^B(X_1)$。对于 $1 < j \le m$，$P_j$ 执行简单哈希 $\mathcal{T}_j^1, \cdots, \mathcal{T}_j^B \gets \mathsf{Simple}_{h_1,h_2,h_3}^B(X_j)$。$P_j$ 进一步定义载荷桶 $\mathcal{V}_j^1, \cdots, \mathcal{V}_j^B$，其中对于 $1 \le b \le B$，$\mathcal{V}_j^b$ 包含 $\mathcal{T}_j^b$ 中元素的关联载荷。
    \item \textbf{带载荷的批量全成员零分享.} 所有参与方调用批大小为 $B$ 的带载荷的批量全成员零分享功能 $\mathcal{F}_{\mathsf{bpMZSp}}$，其中 $P_1$ 作为 $P_\mathsf{pivot}$ 输入 $\mathcal{C}_1^1, \cdots, \mathcal{C}_1^B$，其余每个 $P_j$ 输入 $(\mathcal{T}_j^1, \cdots, \mathcal{T}_j^B)$ 和 $(\mathcal{V}_j^1, \cdots, \mathcal{V}_j^B)$。对于 $1 \le i \le m$，$P_i$ 获得输出份额 $(\textbf{s}_i, \textbf{w}_i)$。
    \item 对于 $1 \le b \le B$，若 $\mathcal{C}_1^b$ 不是空桶，$P_1$ 设置 $w_{1,b} = w_{1,b} + v$，其中 $v$ 是 $\mathcal{C}_1^b$ 中元素的关联载荷。
    \item 对于 $1 \le i \le m$，$P_i$ 分别将 $\textbf{s}_i$ 和 $\textbf{w}_i$ 输入到多方秘密分享洗牌功能 $\mathcal{F}_{\mathsf{shuffle}}$ 中。$P_i$ 接收到打乱后的份额 $\textbf{s}_{i}' = (s_1', \cdots, s_B')$ 和 $\textbf{w}_{i}' = (w_1', \cdots, w_B')$。
    \item 对于 $1 < j  \le m$，$P_j$ 将 $\textbf{s}_j'$ 发送给 $P_1$。$P_1$ 计算重构的指示向量 $\textbf{s}' = \sum_{1 < j  \le m} \textbf{s}_j'$，并定义一个布尔向量 $\textbf{e} = (e_1, \cdots, e_B)$，规则为：若 $s_b' = 0$，则 $e_b = 1$；否则 $e_b = 0$。$P_1$ 将 $\textbf{e}$ 分发给所有 $P_j$。对于 $1 \le i \le m$，$P_i$ 输出向量 $\textbf{e}$ 的汉明重量 $\mathsf{HW}(\textbf{e})$ 作为交集基数。
    \item 对于 $1 \le i \le m$，$P_i$ 在本地计算载荷份额之和 $u_i = \sum_{1 \le b \le B \text{ s.t. } e_b = 1} w_{i,b}'$。对于 $1 < j  \le m$，$P_j$ 将 $u_j$ 发送给 $P_1$。$P_1$ 计算并输出最终的载荷总和 $u = \sum_{1 \le i \le m} u_i$。
\end{enumerate}
\end{trivlist}
\end{minipage}
\end{framed}
\caption{MPSI-card-sum 优化协议}\label{fig:proto-mpsi-sum}
\end{figure}

\section{多方隐私集合求并与并集计算} 
\esection{Multi-Party Private Set Union and Computation on Set Union}

\subsection{多方隐私集合求并、并集求势和基于电路的MPSU}

考虑并集计算 $Y = X_1 \cup \cdots \cup X_m$，其对应的规范集合谓词公式（CSPF）为若干子公式的析取，其中某些子公式直接就是某个原子命题 $x \in X_i$（对于某个 $1 \le i \le m$）。例如，$X_1 \cup \cdots \cup X_m$ 可以表示为 $(x \in X_1) \lor ((x \notin X_1) \land (x \in X_2)) \lor \cdots \lor ((x \notin X_1) \land \cdots \land (x \notin X_{m-1}) \land (x \in X_m))$。在这种情况下，子公式 $x \in X_i$ 仅涉及 $P_i$ 自身，因此 $P_i$ 可以直接在本地对其元素生成秘密分享。特别地，如果 $i = 1$，由于 $P_1$ 是最终接收并集的一方，该子公式甚至可以直接被忽略，只要 $P_1$ 在最终重建的集合中简单地追加自己原本的输入集合 $X_1$ 即可，这避免了大量冗余操作。

\subsection{多方隐私集合求并、并集求势和基于电路的MPSU}

基于上述思路，MPSU 及相关变体的完整优化协议在图 \ref{fig:proto-mpsu} 中给出。

本文的 MPSU 协议遵循了 Liu 和 Gao \cite{LG-ASIACRYPT-2023} 提出的基于秘密分享的 MPSU 范式（该范式的原始设计依赖非共谋假设）。虽然在第三章中，我们通过引入多方秘密分享随机不经意传输（mss-ROT）消除了这一安全缺陷，但 mss-ROT 引入了大量的参与方两两交互（Pairwise Interaction），使得协议的复杂度受到一定影响。

与之形成鲜明对比的是，本章基于通用框架实例化的 MPSU 协议，利用全非成员零分享及其转化技术（结合 Beaver 三元组），成功\textbf{保留了星型通信拓扑（Star Topology）}。这使得本协议不仅能够抵御任意共谋攻击，而且在维持 Leader 为 $O(m^2 n)$ 复杂度的同时，将每个 Client 的计算与通信复杂度压低至 $O(mn)$，首次实现了该基于秘密分享的 MPSU 范式下的最优渐进复杂度。

\begin{figure}[!hbtp]
\begin{framed}
\begin{minipage}[center]{\textwidth}
\begin{trivlist}
\item \textbf{参数：} $m$ 个参与方 $P_1, \cdots, P_m$。集合大小 $n$。元素长度 $l$。全零字符串长度 $l'$。有限域 $\mathbb{F}$。编码函数 $\mathsf{code}:\mathbb{F} \to \{0,1\}^{l+l'}$。布谷鸟哈希参数：哈希函数 $h_1, h_2, h_3$ 及桶数量 $B$。
\item \textbf{输入：} 每个参与方 $P_i$ 持有输入集合 $X_i = \{x_i^1,\cdots, x_i^n\} \subseteq \{0,1\}^l$。

\begin{enumerate}[itemsep=2pt,topsep=0pt,parsep=0pt,leftmargin=15pt]
\item \textbf{哈希分桶.} $P_1$ 执行简单哈希 $\mathcal{T}_1^1, \cdots, \mathcal{T}_1^B \gets \mathsf{Simple}_{h_1,h_2,h_3}^B(X_1)$。对于 $1 < j \le m $，$P_j$ 执行布谷鸟哈希 $\mathcal{C}_j^1, \cdots, \mathcal{C}_j^B \gets \mathsf{Cuckoo}_{h_1,h_2,h_3}^B(X_j)$ 以及简单哈希 $\mathcal{T}_j^1, \cdots, \mathcal{T}_j^B \gets \mathsf{Simple}_{h_1,h_2,h_3}^B(X_j)$。
\item \textbf{批量全非成员零分享.} 对于 $1 < j \le m$，参与方 $P_1, \cdots, P_j$ 调用批大小为 $B$ 的批量全非成员零分享功能 $\mathcal{F}_{\mathsf{bpNMZS}}$。其中 $P_j$ 作为 $P_\mathsf{pivot}$ 输入 $\mathcal{C}_j^1, \cdots, \mathcal{C}_j^B$，其他参与方 $P_{j'}$ ($j' \in \{1,\cdots, j-1\}$) 输入 $\mathcal{T}_{j'}^1, \cdots, \mathcal{T}_{j'}^B$。对于 $1 \le i \le j$，$P_i$ 获得输出份额 $\textbf{s}_{j,i}=(s_{j,i,1},\cdots,s_{j,i,B})$。
\end{enumerate}

后续操作依功能而定：

    \textbf{MPSU.} 
    \begin{enumerate}[itemsep=2pt,topsep=0pt,parsep=0pt,leftmargin=15pt] \addtocounter{enumi}{2}
        \item 对于 $1 < j \le m$ 和 $1 \le b \le B$，若 $\mathcal{C}_j^b$ 不是空桶，$P_j$ 计算 $s_{j,j,b} = \mathsf{code}(s_{j,j,b}) \oplus (x \Vert 0^{l'})$，其中 $x$ 是 $\mathcal{C}_j^b$ 中的元素；否则 $P_j$ 随机采样 $s_{j,j,b} \gets \{0,1\}^{l+l'}$。
        \item 对于 $1 \le i \le m$，$P_i$ 构造一个扩展向量 $\textbf{u}_i \in (\{0,1\}^{l+l'})^{(m-1)B}$，规则如下：对于 $\max(2,i) \le j \le m$ 和 $1 \le b \le B$，$u_{i,(j-2)B+b} = s_{j,i,b}$。将其余位置设为 $0$。
        \item 对于 $1 \le i \le m$，$P_i$ 将 $\textbf{u}_i$ 输入到 $\mathcal{F}_{\mathsf{shuffle}}$ 中，接收打乱后的份额 $\textbf{u}_{i}'$。
        \item 对于 $1 < j \le m$，$P_j$ 将 $\textbf{u}_j'$ 发送给 $P_1$。$P_1$ 计算重构向量 $\textbf{u}' = \sum_{1 < j \le m} \textbf{u}_j'$ 并初始化 $Y = \emptyset$。对于 $1 \le b \le (m-1)B$，如果 $u_b'$ 满足 $y \Vert 0^{l'}$ 的格式（其中 $y \in \{0,1\}^l$），$P_1$ 将 $y$ 加入集合 $Y = Y \cup \{y\}$。最终输出 $Y \cup X_1$。
    \end{enumerate}
    
    \textbf{MPSU-card.} 
    \begin{enumerate}[itemsep=2pt,topsep=0pt,parsep=0pt,leftmargin=15pt] \addtocounter{enumi}{2}
        \item 对于 $1 < j \le m$ 和 $1 \le b \le B$，如果 $\mathcal{C}_j^b$ 是空桶，$P_j$ 随机选择 $s_{j,j,b}$ 的值。
        \item 对于 $1 \le i \le m$，$P_i$ 构造向量 $\textbf{u}_i \in \mathbb{F}^{(m-1)B}$，规则如下：对于 $\max(2,i) \le j \le m$ 和 $1 \le b \le B$，$u_{i,(j-2)B+b} = s_{j,i,b}$。将其余位置设为 $0$。
        \item 对于 $1 \le i \le m$，$P_i$ 将 $\textbf{u}_i$ 输入到 $\mathcal{F}_{\mathsf{shuffle}}$ 中，接收 $\textbf{u}_{i}'$。
        \item 对于 $1 < j \le m$，$P_j$ 将 $\textbf{u}_j'$ 发送给 $P_1$。$P_1$ 输出并集基数为 $n + \mathsf{Zero}(\sum_{1 < j \le m} \textbf{u}_j')$，其中 $\mathsf{Zero}()$ 表示统计重构向量中 0 的数量。
    \end{enumerate}
    
    \textbf{Circuit-MPSU (采用通用电路方案).}
    \begin{enumerate}[itemsep=2pt,topsep=0pt,parsep=0pt,leftmargin=15pt] \addtocounter{enumi}{2}
        \item 对于 $1 < j \le m$ 和 $1 \le b \le B$，若 $\mathcal{C}_j^b$ 不是空桶，$P_j$ 计算 $s_{j,j,b} = \mathsf{code}(s_{j,j,b}) \oplus (x \Vert 0^{l'})$，其中 $x$ 是 $\mathcal{C}_j^b$ 中的元素；否则 $P_j$ 随机采样 $s_{j,j,b} \gets \{0,1\}^{l+l'}$。
        \item 对于 $1 \le i \le m$，$P_i$ 构造扩展向量 $\textbf{u}_i \in (\{0,1\}^{l+l'})^{(m-1)B}$，规则同上，其余位置设为 $0$。
        \item 所有参与方调用包含电路 $C_{(m-1)B,l'}^2$ 的 $m$ 方通用 MPC。对于 $1 \le i \le m, 1 \le k \le (m-1)B$，$P_i$ 将 $u_{i,k}$ 输入到电路中进行联合计算。
    \end{enumerate}
\end{trivlist}
\end{minipage}
\end{framed}
\caption{MPSU / MPSU-card / Circuit-MPSU 优化协议}\label{fig:proto-mpsu}
\end{figure}

\section{各类协议的复杂度分析与对比}
\esection{Implementations and Comprehensive Performance Evaluation}

\section{各类协议的实现与综合性能评估}
\esection{Implementations and Comprehensive Performance Evaluation}

\section{本章小结}
\esection{Summary}